% =====================================================================
%  ch03d.tex —— 3.5 附录：齐次坐标代数的 n-Koszul 性质 + 第 3 章章末评注
%  原书 p.564–567（正文止于 "CHAPTER 4" 之前）
% =====================================================================
% [原文 p.564]
\section{附录：齐次坐标代数的 n-Koszul 性质}

\takeaway{本节补证 §3.3 所用的命题 3.3.4：以齐次坐标代数 $A$ 与理想层复形 $P^\bullet_n$ 为工具，先证明充分大的幂 $\cL=\cM^i$ 满足条件 (43)，再由此推出 $A$ 是 n-Koszul 代数，并得到 $\cL^k$ 的层消解与对角 $\Delta\subset X\times X$ 的 n-消解。}

本附录收集的事实并非原创，它们以各种形式为人熟知。然而，由于缺少一个合适的参考文献，我们不得不按自己所需的形式，给出正文所用断言的证明。这里我们主要使用 [18] 的技巧。

设 $X$ 是光滑射影簇，$\cL$ 是 $X$ 上的极丰富线丛，并满足附加条件 $H^i(X,\cL^k)=0$（对一切 $k>0$ 与 $i\ne 0$）。记 $A$ 为 $X$ 关于 $\cL$ 的齐次坐标代数，即
\begin{equation*}
A=\bigoplus_{k=0}^{\infty}H^0(X,\cL^k).
\end{equation*}

\yizhu{齐次坐标代数（homogeneous coordinate algebra）指射影簇 $X$ 关于极丰富线丛 $\cL$ 的齐次坐标环：把 $\cL$ 各次幂的截面按次数汇成分次代数 $A=\bigoplus_{k\geqslant 0}H^0(X,\cL^k)$。当 $\cL=\cO_X(1)$ 由 $X\hookrightarrow\PP^N$ 给出、且 $X$ 射影正规时，$A$ 就是 $X$ 在 $\PP^N$ 中的齐次坐标环。分次结构使“以 $\cL$ 为权”的上同调能按次数逐块读取，这正是下文把 $A_m=H^0(X,\cL^m)$ 当作系数使用的原因；附加条件 $H^i(X,\cL^k)=0$（$k>0$，$i\ne 0$）则排除了正次数上的高阶上同调。【来源：原书 §3.5 开头及 §3.3 式 (33) 前；Polishchuk--Positselski, \emph{Quadratic Algebras}, AMS Univ. Lecture Ser. 37, 2005, 第 1 章（A 级，经典教材）】}

考虑簇 $X^n$（对某个 $n\in\N$）。下文中记 $\pi^{(n)}_i$ 为 $X^n$ 到第 $i$ 个因子的投影，$\pi^{(n)}_{ij}$ 为它到第 $i$ 与第 $j$ 个因子之积的投影。定义子簇 $\Delta^{(n)}_{(i_1,\ldots,i_k)(i_{k+1},\ldots,i_m)}\subset X^n$ 如下：
\begin{equation*}
\Delta^{(n)}_{(i_1,\ldots,i_k)(i_{k+1},\ldots,i_m)}
:=\bigl\{(x_1,\ldots,x_n)\bigm|x_{i_1}=\cdots=x_{i_k};\ x_{i_{k+1}}=\cdots=x_{i_m}\bigr\}.
\end{equation*}

为简洁起见，我们写 $S^{(n)}_i$ 代替 $\Delta^{(n)}_{(n,n-1,\ldots,i)}$；显然 $S^{(n)}_i\cong X^i$。

现在令
\begin{equation*}
T^{(n)}_i:=\bigcup_{k=1}^{i-1}\Delta^{(n)}_{(n,n-1,\ldots,i)(k,k-1)},
\qquad
\Sigma^{(n)}:=\bigcup_{k=1}^{n}\Delta^{(n)}_{(k,k-1)}.
\end{equation*}

（按定义，$T^{(n)}_1$ 与 $T^{(n)}_2$ 是空子集。）显然 $T^{(n)}_i\subset S^{(n)}_i$。

我们记 $J_{\Sigma^{(n)}}$ 为子概形 $\Sigma^{(n)}\subset X^n$ 的理想层，记 $\cJ^{(n)}_i$ 为 $X^n$ 上的层，它是自然映射 $\cO_{S^{(n)}_i}\to\cO_{T^{(n)}_i}\to 0$ 的核。

% [Tutor Note: 文字层把花体的 $\mathcal{J}$ 一律抽成 $I$（如“I(m) k”）；页图印作 $\mathcal{J}^{(m)}_k$，本译本从页图作 $\cJ$。]

我们暂时固定 $m$ 与 $k\leqslant m$。设 $s$ 是子簇 $S^{(m)}_k\cong X^{k-1}\times X$ 到 $X^n$ 的嵌入，按 $S^{(m)}_k$ 的定义，它在前 $k-1$ 个因子上的作用是恒等，在最后的第 $k$ 个因子上的作用是对角。记 $p$ 为 $S^{(m)}_k$ 到 $X^{k-1}$ 的投影，即前 $k-1$ 个因子之积。

% [Tutor Note: 页图与文字层此处均印作嵌入到 $X^n$；但 $s$ 应是 $S^{(m)}_k\subset X^m$ 到 $X^m$ 的嵌入（引理 3.5.1 的各项都在 $X^m$ 上取上同调），据上下文应为 $X^m$。此处照录。]

\begin{lemma}{3.5.1}
层 $\cJ^{(m)}_1$ 同构于 $\cO_{\Delta^{(m)}_{(m,\ldots,1)}}$；而对 $k>1$，$\cJ^{(m)}_k$ 同构于 $s_*p^*(J_{\Sigma^{(k-1)}})$。特别地，当 $k>1$ 时有同构
\begin{enumerate}[label=\alph*),leftmargin=2.8em,itemsep=2pt,topsep=2pt]
\item $H^j\bigl(X^m,\cJ^{(m)}_k\otimes(\cL\boxtimes\cdots\boxtimes\cL^i)\bigr)
=H^j\bigl(X^{k-1},J_{\Sigma^{(k-1)}}\otimes(\cL\boxtimes\cdots\boxtimes\cL)\bigr)\otimes A_{m-k+i}$，对一切 $i>0$；
\item $R^j\pi^{(m)}_{1*}\bigl(\cJ^{(m)}_k\otimes(\cO\boxtimes\cL\boxtimes\cdots\boxtimes\cL^i)\bigr)
\cong R^j\pi^{(k-1)}_{1*}\bigl(J_{\Sigma^{(k-1)}}\otimes(\cO\boxtimes\cL\boxtimes\cdots\boxtimes\cL)\bigr)\otimes A_{m-k+i}$，对一切 $i>0$；
\item $R^j\pi^{(m)}_{1m*}\bigl(\cJ^{(m)}_k\otimes(\cO\boxtimes\cL\boxtimes\cdots\boxtimes\cL^i)\bigr)
\cong R^j\pi^{(k-1)}_{1*}\bigl(J_{\Sigma^{(k-1)}}\otimes(\cO\boxtimes\cL\boxtimes\cdots\boxtimes\cL)\bigr)\boxtimes\cL^{m-k+i}$，对一切 $i$。
\end{enumerate}
\end{lemma}

\begin{proof}
对 $k>1$，$\cJ^{(m)}_k$ 同构于 $s_*p^*(J_{\Sigma^{(k-1)}})$ 这一断言，直接由 $\cJ^{(m)}_k$ 以及子概形 $T^{(m)}_k$、$S^{(m)}_k$ 的定义得到。其余断言由它直接推出。
\end{proof}

对 $n$ 作归纳，容易看出复形
\begin{equation*}
P^\bullet_n:\ 0\to J_{\Sigma^{(n)}}\to\cJ^{(n)}_n\to\cJ^{(n)}_{n-1}\to\cdots\to\cJ^{(n)}_2\to\cJ^{(n)}_1\to 0
\end{equation*}
% [原文 p.565]
在 $X^n$ 上正合。例如，当 $n=2$ 时，这个复形就是 $X\times X$ 上的短正合列
\begin{equation*}
P^\bullet_2:\ 0\to J_\Delta\to\cO_{X\times X}\to\cO_\Delta\to 0.
\end{equation*}

\yizhu{Koszul 复形（Koszul complex）原指由环中元素序列 $f_1,\dots,f_n$ 以 $\bigwedge^\bullet$ 为项作出的复形：它的正合性等价于该序列构成正则序列，这是“线性消解”的原型。本节的 $P^\bullet_n$ 取同样的形状，项换成理想层 $J_{\Sigma^{(n)}}$ 与 $\cJ^{(n)}_i$，整个复形在 $X^n$ 上正合，给出 $J_{\Sigma^{(n)}}$ 的一个消解；$n=2$ 时它退化为 $0\to J_\Delta\to\cO_{X\times X}\to\cO_\Delta\to 0$。附录正是借它把“对角有消解”与“$A$ 是 n-Koszul 代数”两件事接起来。【来源：原书 §3.5；Koszul 复形的经典表述见 Weibel, \emph{An Introduction to Homological Algebra}, CUP 1994, §4.5（A 级）】}

\begin{lemma}{3.5.2}
设 $X$ 是带丰富线丛 $\cM$ 的光滑射影簇。则对任一正整数 $k$，存在 $i$，使得丛 $\cL=\cM^i$ 对一切 $1<m\leqslant k$ 具有如下性质：
\begin{equation}
\left.
\begin{aligned}
\text{a)}\quad & H^j\bigl(X^m,J_{\Sigma^{(m)}}\otimes(\cL\boxtimes\cdots\boxtimes\cL)\bigr)=0
&&\text{对 } j\ne 0;\\
\text{b)}\quad & R^j\pi^{(m)}_{1*}\bigl(J_{\Sigma^{(m)}}\otimes(\cO\boxtimes\cL\boxtimes\cdots\boxtimes\cL)\bigr)=0
&&\text{对 } j\ne 0;\\
\text{c)}\quad & R^j\pi^{(m)}_{1m*}\bigl(J_{\Sigma^{(m)}}\otimes(\cO\boxtimes\cL\boxtimes\cdots\boxtimes\cL\boxtimes\cO)\bigr)=0
&&\text{对 } j\ne 0.
\end{aligned}
\right\}\tag{43}
\end{equation}
\end{lemma}

\begin{proof}
对任一 $m$，$X^m$ 上的线丛 $\cM\boxtimes\cdots\boxtimes\cM$、$\cO\boxtimes\cM\boxtimes\cdots\boxtimes\cM$ 与 $\cO\boxtimes\cM\boxtimes\cdots\boxtimes\cM\boxtimes\cO$ 分别是丰富丛、$\pi^{(m)}_1$-丰富丛与 $\pi^{(m)}_{1m}$-丰富丛。因此，对它们中的每一个，都存在一个整数，使得性质 a)、b)、c) 对这几个丛的一切高于该整数的幂成立。取这些整数在一切 $m\leqslant k$ 上的最大值，并记为 $i$。于是性质 a)、b)、c) 对 $\cL=\cM^i$ 也成立。
\end{proof}

我们引入如下记号：
\begin{equation*}
\begin{aligned}
&B_n:=H^0\bigl(X^n,J_{\Sigma^{(n)}}\otimes(\cL\boxtimes\cdots\boxtimes\cL)\bigr)\\
&\text{与}\quad R_{n-1}:=R^0\pi^{(n)}_{1*}\bigl(J_{\Sigma^{(n)}}\otimes(\cO\boxtimes\cL\boxtimes\cdots\boxtimes\cL)\bigr).
\end{aligned}
\end{equation*}

\begin{proposition}{3.5.3}
设 $\cL$ 是光滑射影簇 $X$ 上的极丰富丛，且对一切满足 $1<m\leqslant n+\dim X+2$ 的 $m$ 满足条件 (43)。则
\begin{enumerate}[label=\arabic*),leftmargin=2.6em,itemsep=2pt,topsep=2pt]
\item $A$ 是 n-Koszul 代数，即序列
\begin{equation*}
B_n\otimes_k A\to B_{n-1}\otimes_k A\to\cdots\to B_1\otimes_k A\to A\to k\to 0
\end{equation*}
正合；
\item $X$ 上如下形式的层复形
\begin{equation*}
\begin{aligned}
&A_{k-n}\otimes R_n\to A_{k-n+1}\otimes R_{n-1}\to\\
&\qquad\cdots\to A_{k-1}\otimes R_1\to A_k\otimes R_0\to\cL^k\to 0
\end{aligned}
\end{equation*}
对任一 $k\geqslant 0$ 正合（若 $k-i<0$，则按定义 $A_{k-i}=0$）；
\item 层复形
\begin{equation*}
\cL^{-n}\boxtimes R_n\to\cdots\to\cL^{-1}\boxtimes R_1\to\cO_M\boxtimes R_0\to\cO_\Delta
\end{equation*}
在 $X\times X$ 上正合；也就是说，它给出 $X\times X$ 的对角的一个 n-消解。
\end{enumerate}
\end{proposition}

\yizhu{n-Koszul 性质的条件已由 §3.3 定义 3.3.3 给出：$k$ 作为右 $A$-模有自由消解 $B_n\otimes_k A\to\cdots\to B_1\otimes_k A\to A\to k\to 0$，其中 $B_1=A_1$，$B_m=\Ker(B_{m-1}\otimes A_1\to B_{m-2}\otimes A_2)$ 逐次取核而得（式 (33)）。$n=2$ 时这就是经典的 Koszul 代数，一般 $n$ 只是把“线性消解”截到第 $n$ 步。判断 Koszul 性的标准不变量是 Ext 代数 $\Ext^\bullet_A(k,k)=\bigoplus_i\Ext^i_A(k,k)$：$A$ 为 Koszul 代数，当且仅当这一分次代数由 $\Ext^1_A(k,k)$ 生成。附录不去计算 Ext 代数，而是直接给出使上述消解正合的充分条件 (43)。【来源：原书 §3.3 定义 3.3.3 与式 (33)、§3.5 命题 3.5.3；Polishchuk--Positselski, \emph{Quadratic Algebras}, AMS Univ. Lecture Ser. 37, 2005, 第 1--2 章（A 级，经典教材）】}

% [Tutor Note: 本命题即 §3.3 命题 3.3.4 的“证明版本”：把那里的簇 $M$ 改写为 $X$，并显式给出条件 (43)。第 3) 条中的 $\cO_M$ 沿用了 §3.3 的记号——那里规定 $R_0=\cO_M$；本节簇改称 $X$，故 $\cO_M$ 与 $\cO_X$ 是同一个层，并非笔误。此处照录。]

\yizhu{对角的消解（resolution of the diagonal）指 $X\times X$ 上以结构层 $\cO_\Delta$ 为末项的复形。它的用处在于：把 $\cO_\Delta$ 换成由外张量积 $\cL^a\boxtimes\cL^b$ 这类简单对象组成的复形之后，关于“函子的核”的论断就可以逐项验证。命题 3.5.3 第 3) 条给出的正是一条长度 $n$ 的这类消解 $\cL^{-n}\boxtimes R_n\to\cdots\to\cO_M\boxtimes R_0\to\cO_\Delta$（§3.3 规定 $R_0=\cO_M$），它与 $A$ 的 n-Koszul 消解互为表里；§3.3 在 $\PP^N\times\PP^N$ 上先用的是 Beilinson 型消解（文献 [2]），附录补上 $X\times X$ 上与之配套的 n-消解。【来源：原书 §3.3 命题 3.3.4 及其后的用法、§3.5 命题 3.5.3 与式 (46)】}

% [原文 p.566]
\begin{proof}
\textbf{1)} 首先，把引理 3.5.1 与引理 3.5.2 结合起来，对任一 $1<m\leqslant n+\dim X+2$，我们看到
\begin{align}
\text{1)}\quad & H^0\bigl(X^m,\cJ^{(m)}_k\otimes(\cL\boxtimes\cdots\boxtimes\cL^i)\bigr)\notag\\
&\qquad =H^0\bigl(X^{k-1},J_{\Sigma^{(k-1)}}\otimes(\cL\boxtimes\cdots\boxtimes\cL)\bigr)\otimes A_{m-k+i}
=B_{k-1}\otimes A_{m-k+i};\notag\\
\text{2)}\quad & H^j\bigl(X^m,\cJ^{(m)}_k\otimes(\cL\boxtimes\cdots\boxtimes\cL^i)\bigr)=0
\qquad\text{对 } j\ne 0.\tag{44}
\end{align}

考虑复形 $P^\bullet_m\otimes(\cL\boxtimes\cdots\boxtimes\cL)$（$m\leqslant n+\dim X+1$）。对它们施加 $H^0$，并用条件 (44)，我们得到正合列
\begin{equation*}
0\to B_m\to B_{m-1}\otimes_k A_1\to\cdots\to B_1\otimes_k A_{m-1}\to A_m\to 0
\end{equation*}
对任一 $m\leqslant n+\dim X+1$ 成立。

现在令 $m=n+\dim X+1$，并记 $W^\bullet_m$ 为复形
\begin{equation*}
\cJ^{(m)}_m\to\cJ^{(m)}_{m-1}\to\cdots\to\cJ^{(m)}_2\to\cJ^{(m)}_1\to 0,
\end{equation*}
它是 $J_{\Sigma^{(m)}}$ 的一个右消解。我们取复形
\begin{equation*}
W^\bullet_m\otimes(\cL\boxtimes\cdots\boxtimes\cL\boxtimes\cL^i)
\end{equation*}
并对它施加函子 $H^0$。我们得到序列
\begin{equation*}
B_{m-1}\otimes_k A_i\to B_{m-2}\otimes_k A_{i+1}\to\cdots\to B_1\otimes_k A_{m+i-2}\to A_{m+i-1}\to 0.
\end{equation*}
由 (44) 2) 可知，这个复形的上同调等于
\begin{equation*}
H^j\bigl(X^m,J_{\Sigma^{(m)}}\otimes(\cL\boxtimes\cdots\boxtimes\cL\boxtimes\cL^i)\bigr).
\end{equation*}
而 (43) b) 给出
\begin{equation*}
\begin{aligned}
&H^j\bigl(X^m,J_{\Sigma^{(m)}}\otimes(\cL\boxtimes\cdots\boxtimes\cL\boxtimes\cL^i)\bigr)\\
&\qquad =H^j\bigl(X,R^0\pi^{(m)}_{m*}(J_{\Sigma^{(m)}}\otimes(\cL\boxtimes\cdots\boxtimes\cL\boxtimes\cO))\otimes\cL^i\bigr).
\end{aligned}
\end{equation*}
因此，当 $j>\dim X$ 时该上同调为零，于是有形式如下的正合列
\begin{equation*}
B_n\otimes_k A_{m-n+i-1}\to B_{n-1}\otimes_k A_{m-n+i}\to\cdots\to B_1\otimes_k A_{m+i-2}\to A_{m+i-1}\to 0
\end{equation*}
（$i\geqslant 1$）。然而，$i\leqslant 1$ 情形的正合性上面已经证明。于是 $A$ 是 n-Koszul 代数。

% [Tutor Note: 页图印作 “the exactness for $i\leqslant 1$”，而上一句的序列是对 $i\geqslant 1$ 给出的；按上下文这里指 $i=1$ 的情形。此处照录。]

\textbf{2)} 证明与 1) 类似。我们有同构
\begin{align}
\text{1)}\quad & R^0\pi^{(m)}_{1*}\bigl(\cJ^{(m)}_k\otimes(\cO\boxtimes\cL\boxtimes\cdots\boxtimes\cL)\bigr)\notag\\
&\qquad\cong R^0\pi^{(k-1)}_{1*}\bigl(J_{\Sigma^{(k-1)}}\otimes(\cO\boxtimes\cL\boxtimes\cdots\boxtimes\cL)\bigr)\otimes A_{m-k+1};\notag\\
\text{2)}\quad & R^j\pi^{(m)}_{1*}\bigl(\cJ^{(m)}_k\otimes(\cO\boxtimes\cL\boxtimes\cdots\boxtimes\cL)\bigr)=0
\qquad\text{对一切 } j\ne 0.\tag{45}
\end{align}

% [原文 p.567]
对复形 $P^\bullet_m\otimes(\cO\boxtimes\cL\boxtimes\cdots\boxtimes\cL)$（$m\leqslant n+\dim X+2$）施加函子 $R^0\pi^{(m)}_{1*}$，我们得到 $X$ 上一个正合的层复形
\begin{equation*}
0\to R_{m-1}\to A_1\otimes R_{m-2}\to\cdots\to A_{m-2}\otimes R_1\to A_{m-1}\otimes R_0\to\cL^{m-1}\to 0
\end{equation*}
（$m\leqslant n+\dim X+2$）。

我们考虑 $m=n+\dim X+2$ 的情形。对
\begin{equation*}
W^\bullet_m\otimes(\cO\boxtimes\cL\boxtimes\cdots\boxtimes\cL\boxtimes\cL^i)
\end{equation*}
施加函子 $R^0\pi^{(m)}_{1*}$，得到复形
\begin{equation*}
A_i\otimes R_{m-2}\to\cdots\to A_{m+i-3}\otimes R_1\to A_{m+i-2}\otimes R_0\to\cL^{m+i-2}\to 0.
\end{equation*}
由性质 (45)，它的上同调为
\begin{equation*}
\begin{aligned}
&R^j\pi^{(m)}_{1*}\bigl(J_{\Sigma^{(m)}}\otimes(\cO\boxtimes\cL\boxtimes\cdots\boxtimes\cL\boxtimes\cL^i)\bigr)\\
&\qquad\cong R^jp_{1*}\bigl(R^0\pi^{(m)}_{1m*}(J_{\Sigma^{(m)}}\otimes(\cO\boxtimes\cL\boxtimes\cdots\boxtimes\cL\boxtimes\cO))\otimes(\cO\boxtimes\cL^i)\bigr),
\end{aligned}
\end{equation*}
它在 $j>\dim X$ 时为零。于是 $X$ 上的层序列
\begin{equation*}
A_{k-n}\otimes R_n\to A_{k-n+1}\otimes R_{n-1}\to\cdots\to A_{k-1}\otimes R_1\to A_k\otimes R_0\to\cL^k\to 0
\end{equation*}
对一切 $k\geqslant 0$ 正合。

\textbf{3)} 考虑 $X^{n+2}$ 上的复形 $W^\bullet_{n+2}\otimes(\cO\boxtimes\cL\boxtimes\cdots\boxtimes\cL\boxtimes\cL^{-n})$。对它施加 $R^0\pi^{(n+2)}_{1(n+2)*}$，我们得到 $X\times X$ 上如下复形：
\begin{equation}
\cL^{-n}\boxtimes R_n\to\cdots\to\cL^{-1}\boxtimes R_1\to\cO_M\boxtimes R_0\to\cO_\Delta.\tag{46}
\end{equation}

回忆引理 3.5.1 的条件 c) 与条件 (43) b)，我们看到它的上同调层同构于
\begin{equation*}
R^j\pi^{(n+2)}_{1(n+2)*}\bigl(J_{\Sigma^{(n+2)}}\otimes(\cO\boxtimes\cL\boxtimes\cdots\boxtimes\cL\boxtimes\cO)\bigr)\otimes(\cO\boxtimes\cL^{-n}),
\end{equation*}
而由 (43) c)，它在 $j>0$ 时消没。也就是说，复形 (46) 正合。
\end{proof}

% [Tutor Note: 式 (46) 与命题 3.5.3 第 3) 条一样沿用 $\cO_M$（即本节语境下的 $\cO_X$），见前注；此处照录。]

\chapterreview{%
第 3 章是全文的核心。引言把全书归结为两个问题：两个簇的凝聚层导出范畴何时等价，给定簇的正合自等价群是什么；而回答它们的前提，是先说清范畴之间的函子从何而来。本章证明：光滑射影簇之间的正合函子只要全忠实并且有伴随，就同构于某个由积上的对象定义的函子 $\Phi_E$（定理 3.2.1）；等价作为全忠实、且以拟逆为伴随的函子自然满足条件，于是每个等价都由 $M\times X$ 上同构唯一的对象 $E$ 表示（定理 3.2.2）。这条“函子必由核对象表示”的定理，是第 2 章那批具体等价（flop）与第 4、5 章全部分类陈述的公共入口。

全章的技术分工十分清楚。3.1 处理三角范畴内部的粘合：Postnikov 图与卷积把一族逐次扩张的三角拼成一个对象，并给出卷积存在、且在同构意义下（乃至典范地）唯一的充分条件，这是后文把局部构造拼成整体对象的语言。3.2 陈述主定理，并先解决有界性：引理 3.2.3 说明有左伴随的正合函子必有界，注 3.2.4 由此把 $F$ 的像的上同调压进一个固定区间 $[-a,0]$，归纳才得以进行。3.3 是构造：先把 $M$ 闭嵌入 $\PP^N$，在 $\PP^N\times X$ 上造出 $E'$，再借命题 3.3.5 把它下降为 $M\times X$ 上的对象 $E$；这一步真正用到了齐次坐标代数的 n-Koszul 性质与对角的 n-消解（命题 3.3.4）。3.4 是验证：在阿贝尔范畴中引入丰富序列（ample sequence），用引理 3.4.3--3.4.5 与命题 3.4.6，把在丰富序列上给出的自然变换延拓到整个导出范畴，从而证明 $F\cong\Phi_E$。3.5 是附录，为 3.3 所用的命题 3.3.4 补一个自足的证明。

核心结论可以一句话概括：等价不是“碰巧”由某个层或复形给出，而是必然如此。证明的路线是先构造、后验证——先在射影空间里造出候选核，再证明它下降，并证明由它定义的函子与原函子在每个对象上一致。这条路线之所以走得通，关键在于导出范畴的函子由它在生成元（这里是线丛族 $\cL^k$）上的值决定；两个函子之间的自然变换一旦在 ample sequence 上相等，就被迫整体相等。

辨析与提醒有三点。其一，“全忠实且有伴随”这个条件不能去掉：一般正合函子是否也能由积上的对象表示，正是 §3.2 开头记为未解的那个猜想。其二，唯一性只是同构意义下的，核对象 $E$ 还可以相差平移与线丛 twist，第 5 章描述自等价群时利用的正是这一点。其三，3.5 的 n-Koszul 论断并非作者原创，作者也明说这些事实“以各种形式为人熟知”；它在此处的地位是技术引理，读者不必把它当作新的理论贡献，但其中“由条件 (43) 推出 n-Koszul”与“对角有 n-消解”两步，是把代数性质翻译成层论消解的范例。

从更长的视野看，本章完成了一次视角转换：不再把导出范畴的等价看作两个范畴之间的抽象同构，而是看作积上的一个对象。于是“等价”成了可计算、可分类的几何数据，第 4 章关于 K3 曲面导出范畴的分类、第 5 章对自等价群的描述，都是在这一转换之后才成为可能的陈述。也正因如此，3.5 这样一段纯技术的附录并不多余：没有对角消解与 n-Koszul 性质，3.3 的构造就缺一块砖，整章“函子—核对象”的对应也无法闭合。}
